\documentclass[preview]{standalone}

\usepackage{amsmath}
\usepackage{amssymb}
\usepackage{stellar}
\usepackage{definitions}
\usepackage{bettelini}

\begin{document}

\id{graph-eulerian-walks}
\genpage

\section{Eulerian walks}

\begin{snippetdefinition}{eulerian-circuit-definition}{Eulerian circuit}
    Let \(G\) be a finite connected \graph.
    An \emph{Eulerian circuit} of \(G\) is a closed \walk, that is, a walk
    whose origin and terminal vertex coincide, that traverses every edge of
    \(G\) exactly once.
    Vertices are allowed to repeat.
\end{snippetdefinition}

\begin{snippet}{konigsberg-eulerian-interpretation}
    The Königsberg bridges problem asks whether the graph associated with the
    bridges admits an Eulerian circuit.
\end{snippet}

\begin{snippettheorem}{graph-eulerian-circuit-theorem}{Euler's theorem for Eulerian circuits}
    Let \(G\) be a finite connected \graph. Then \(G\) admits an Eulerian
    circuit \ifandonlyif every vertex of \(G\) has even degree.
\end{snippettheorem}

\begin{snippetproof}{graph-eulerian-circuit-theorem-proof}{graph-eulerian-circuit-theorem}{Euler's theorem for Eulerian circuits}
    \begin{iffproof}
        \item Suppose that \(G\) admits an Eulerian circuit
        \[
            v_0,e_1,v_1,e_2,\ldots,e_n,v_n,
            \qquad v_n=v_0.
        \]
        Every time the walk passes through a vertex, it uses one edge to enter
        and one edge to leave. Thus the edges incident with each vertex are
        paired. The initial vertex satisfies the same condition: one edge is
        used to leave it at the beginning and one edge is used to enter it at
        the end. Hence every vertex has even degree.

        \item Suppose that every vertex of \(G\) has even degree.
        Start at any vertex \(v_0\) and traverse edges that have not yet been
        used, choosing an available unused edge at each step. Since \(G\) is
        finite, the process eventually stops.

        It can stop only at \(v_0\). Indeed, if it arrived at a vertex
        different from \(v_0\), it would have used one edge to enter that
        vertex. Since the degree is even, the used incident edges at that
        vertex can be paired as entries and exits, so at least one unused
        incident edge would remain available. Therefore the walk cannot get
        stuck away from \(v_0\).

        We have constructed a closed circuit \(C\). If \(C\) contains every
        edge of \(G\), then \(C\) is Eulerian. Otherwise, because \(G\) is
        connected, some vertex of \(C\) is incident with an unused edge.
        Starting from that vertex and using only unused edges, the same
        argument constructs another closed circuit. Insert this new circuit
        into \(C\), obtaining a longer closed circuit.

        Since \(G\) is finite, repeating this process eventually produces a
        closed walk containing every edge exactly once. This is an Eulerian
        circuit.
    \end{iffproof}
\end{snippetproof}

\begin{snippetcorollary}{konigsberg-no-eulerian-circuit-corollary}{Königsberg bridges}
    The graph of the Königsberg bridges does not admit an Eulerian circuit.
\end{snippetcorollary}

\begin{snippetproof}{konigsberg-no-eulerian-circuit-proof}{konigsberg-no-eulerian-circuit-corollary}{Königsberg bridges}
    In the graph of the Königsberg bridges, the four regions have degrees
    \[
        3,\ 3,\ 3,\ 5.
    \]
    All four degrees are odd. By
    \snippetref[graph-eulerian-circuit-theorem][Euler's theorem for Eulerian circuits],
    no Eulerian circuit exists.
\end{snippetproof}

\begin{snippetdefinition}{open-eulerian-trail-definition}{Open Eulerian trail}
    Let \(G\) be a finite connected \graph.
    An \emph{open Eulerian trail} of \(G\) is a \walk that traverses every edge
    of \(G\) exactly once and whose origin and terminal vertex are distinct.
\end{snippetdefinition}

\begin{snippettheorem}{open-eulerian-trail-theorem}{Open Eulerian trail criterion}
    Let \(G\) be a finite connected \graph.
    There exists an open Eulerian trail in \(G\) \ifandonlyif \(G\) has exactly
    two vertices of odd degree.
    In that case, the two vertices of odd degree are necessarily the origin and
    terminal vertex of the trail.
\end{snippettheorem}

\begin{snippetproof}{open-eulerian-trail-theorem-proof}{open-eulerian-trail-theorem}{Open Eulerian trail criterion}
    Suppose that an open Eulerian trail exists from \(u\) to \(v\), with
    \(u\neq v\). Every intermediate vertex is entered and left the same number
    of times, so it has even degree. The vertex \(u\) has one more exit than
    entrance, and \(v\) has one more entrance than exit. Hence \(u\) and \(v\)
    have odd degree, and every other vertex has even degree.

    Conversely, suppose that \(G\) has exactly two vertices of odd degree,
    say \(u\) and \(v\). Add a new edge \(e\) joining \(u\) and \(v\). In the
    new graph, every vertex has even degree. By
    \snippetref[graph-eulerian-circuit-theorem][Euler's theorem for Eulerian circuits],
    the new graph admits an Eulerian circuit. Removing the added edge \(e\)
    from that circuit leaves a walk using every original edge exactly once,
    with origin \(u\) and terminal vertex \(v\). Thus \(G\) admits an open
    Eulerian trail.
\end{snippetproof}

\end{document}
